\documentclass[11pt,a4paper]{article}
\usepackage[T1]{fontenc}
\usepackage[utf8]{inputenc}
\usepackage{amsmath,amssymb,graphicx,hyperref,longtable,booktabs,array}
\usepackage[margin=2.2cm]{geometry}
\providecommand{\tightlist}{\setlength{\itemsep}{0pt}\setlength{\parskip}{0pt}}
\setlength{\parindent}{0pt}\setlength{\parskip}{6pt}

\begin{document}
\hypertarget{piece-8-ablations-posterior-coverage-and-final-report-26-nov-2026}{%
\section{Piece 8: ablations, posterior coverage and final report (26 Nov
2026)}\label{piece-8-ablations-posterior-coverage-and-final-report-26-nov-2026}}

Code: \texttt{code/run\_ablate.py}, \texttt{code/run\_coverage.py},
\texttt{code/inputs.py} (AR(1) noise state and posterior sd added; tests
9 of 9 pass); results \texttt{results/ablate\_results.json},
\texttt{results/coverage\_results.json}, figure
\texttt{results/fig\_pieces678.png}; report
\texttt{docs/tex/final-report.tex}. Written by Claude (AI assistant);
the team must read and sign off before claiming it.

\hypertarget{ablations-rapid-change-rmse-in-mgdl-median-over-27-patients-days-4-to-5-tau--10-s0s1-and-noise-28-difference--full-model-minus-variant-95-ci}{%
\subsection{Ablations (rapid-change RMSE in mg/dL, median over 27
patients, days 4 to 5, tau = 10; S0/S1 and noise 2/8; difference = full
model minus variant, 95\%
CI)}\label{ablations-rapid-change-rmse-in-mgdl-median-over-27-patients-days-4-to-5-tau--10-s0s1-and-noise-28-difference--full-model-minus-variant-95-ci}}

\begin{longtable}[]{@{}lll@{}}
\toprule\noalign{}
Variant & filter (S0-2, S0-8, S1-2, S1-8) & smoother (S0-2, S0-8, S1-2,
S1-8) \\
\midrule\noalign{}
\endhead
\bottomrule\noalign{}
\endlastfoot
no inputs (random walk) & 7.8, 20.3, 16.0, 25.8 & 4.9, 8.3, 13.5,
13.7 \\
full input model & 8.4, 14.9, 18.9, 21.2 & 4.6, 8.4, 11.9, 12.1 \\
gains frozen at prior means & 11.1, 19.7, 18.1, 25.7 & 4.2, 8.0, 11.3,
11.9 \\
kernel times x 0.5 & 8.2, 15.8, 14.7, 20.2 & 4.5, 8.3, 12.4, 12.3 \\
kernel times x 2 & 7.4, 13.3, 16.3, 19.4 & 4.6, 8.0, 13.3, 14.2 \\
+ AR(1) noise state & 8.9, 14.0, 17.7, 21.3 & 4.9, 8.0, 11.9, 11.9 \\
\end{longtable}

\begin{itemize}
\tightlist
\item
  Gain learning matters for the real-time filter (frozen gains cost +2.7
  to +4.8 mg/dL, as differences of medians, in three of four conditions)
  but not for the smoother (differences +0.13 {[}+0.10, +0.35{]}, +0.20
  {[}-0.43, +0.62{]}, -0.02 {[}-0.33, +0.53{]}, +1.09 {[}-0.43,
  +1.47{]}).
\item
  Kernel shape is forgiving: halving the time constants changes the
  smoother by -0.1 to +0.7 mg/dL (variant minus full model); doubling
  them costs 0.1 to 1.1 mg/dL in all four conditions (full model minus
  variant -0.08, -1.11, -0.86, -0.95).
\item
  The AR(1) noise state does not improve point accuracy: smoother
  difference to the full model -0.00 {[}-0.03, +0.22{]}, -0.22 {[}-0.37,
  +0.24{]}, +0.26 {[}+0.08, +0.40{]}, -0.19 {[}-0.45, 0.00{]}; real-time
  filter -0.18 to -0.71 (CI excludes 0 in two conditions).
\end{itemize}

\hypertarget{posterior-coverage-95-interval-for-blood-glucose-smoother-fraction-of-days-4-to-5-inside-median-over-patients-noise-settings-chosen-per-patient-by-likelihood-on-days-1-to-3}{%
\subsection{Posterior coverage (95\% interval for blood glucose,
smoother, fraction of days 4 to 5 inside; median over patients; noise
settings chosen per patient by likelihood on days 1 to
3)}\label{posterior-coverage-95-interval-for-blood-glucose-smoother-fraction-of-days-4-to-5-inside-median-over-patients-noise-settings-chosen-per-patient-by-likelihood-on-days-1-to-3}}

\begin{longtable}[]{@{}lllll@{}}
\toprule\noalign{}
Condition & white noise, no inputs & white noise + inputs & AR(1) noise
+ inputs & AR(1) noise, no inputs \\
\midrule\noalign{}
\endhead
\bottomrule\noalign{}
\endlastfoot
S0, noise 2 & 0.70 & 0.67 & 0.89 & 0.91 \\
S0, noise 8 & 0.66 & 0.59 & 0.89 & 0.91 \\
S1, noise 2 & 0.52 & 0.50 & 0.70 & 0.74 \\
S1, noise 8 & 0.59 & 0.58 & 0.88 & 0.87 \\
\end{longtable}

At rapid-change times the AR(1) model covers 0.57, 0.82, 0.08 and 0.59
(S0-2, S0-8, S1-2, S1-8) against 0.34, 0.49, 0.04 and 0.28 for white
noise with inputs. When noise settings are chosen for RMSE instead of
likelihood, smoother coverage is 0.86, 0.17, 0.42, 0.15: RMSE-tuned
settings give overconfident intervals. The AR structure used is the
simulator\textquotesingle s own (30 min, 60\% of the variance), so this
is an oracle-structure ablation; real sensor noise structure is unknown.

\hypertarget{conclusions}{%
\subsection{Conclusions}\label{conclusions}}

\begin{enumerate}
\def\labelenumi{\arabic{enumi}.}
\tightlist
\item
  The AR(1) noise state is for uncertainty, not accuracy: it brings 95\%
  intervals to about 0.9 coverage in the matched sensor case but does
  not lower RMSE.
\item
  When the sensor adds lag the model does not know (S1 at noise 2),
  intervals fail at rapid change (0.04 to 0.08 coverage): model
  uncertainty from the lag cannot be fixed by noise modelling.
\item
  See \texttt{docs/tex/final-report.tex} for the integrated account.
\end{enumerate}
\end{document}
